\section{Details of the marked spectral-curve classification}
\label{sec:geometry-details}

We prove the marked-curve statements of \cref{sec:spectral-geometry},
including existence and reconstruction, to justify the fixed-size
classification used in the spectral discussion.

\subsection{Existence of regular spectral curves}

The regular locus is nonempty for every positive $L,M,N$, with
phases allowed to vary. To see this, choose intercepts $e_h$ distinct
within each group so that the lines
$\lambda-e_h-s_h\nu=0$, $1\leq h\leq d$, have no triple affine
intersection. Set $X=\diag(e_h)+\varepsilon A$, with zero diagonal
blocks in $A$ and $A_{hi}A_{ih}\ne0$ between different groups.
At $\varepsilon=0$ the curve is the union of these lines.
Near an intersection of lines $h,i$ of different slopes, use their
equations as local coordinates $x,y$. A Schur complement gives,
up to a nonvanishing factor,
\[
 xy-\varepsilon^2 A_{hi}A_{ih}
   +O\bigl(\varepsilon^2x,\varepsilon^2y,\varepsilon^3\bigr).
\]
Its critical point has $x,y=O(\varepsilon^2)$, where the displayed
nonzero quadratic term prevents a singularity for small nonzero
$\varepsilon$. Blowing up the three infinity points separates their
distinct tangent branches. They meet the exceptional curves
transversely, so remain smooth there uniformly for small $\varepsilon$.
Compactness then
excludes further singularities. On this blown-up surface the central
fiber is a connected reduced union of lines, so nearby smooth fibers
are connected, hence irreducible. Thus small generic couplings give
regular curves. Smoothness and irreducibility are open conditions,
so they hold on a nonempty Zariski-open subset of the parameter space.

\subsection{Proof of the marked-curve characterization}

\begin{proof}[Proof of \cref{thm:geometry-classification}]
The implication from~1 to~2 follows from
\cref{prop:geometry-genus}. The expansions
\cref{eq:geometry-branches} give~3: subtracting or adding $\nu$
cancels the pole on the indicated branch family. The two affine
coordinates generate the function field because the plane model
is birational to $C$.

Conversely, the pole conditions in~3 imply that $\nu$ has exactly
$d$ simple poles, one at each marked point. The functions
$1,\lambda,\nu$ define a map to the plane with no base points:
at each marked point at least one of $\lambda,\nu$ has a simple
pole, and away from them the constant section is nonzero.
The map is birational because $\lambda,\nu$ generate $\C(C)$.
The image has degree $d$, since a generic linear combination of
these three functions has $d$ poles. The three pole cancellations
put the marked branches at $[1:1:0]$, $[-1:1:0]$, $[0:1:0]$.
Their simple poles make the branches smooth and transverse to infinity.
The genus assumption forces their total singularity defect to be
exactly the sum in \cref{prop:geometry-genus}. Hence their tangents
are distinct and there are no additional singularities, proving~2.

Finally, a plane model in~2 can be reconstructed as a pencil
$\mathsf X+\nu\mathsf D$. The spectral line-bundle construction
\cite{Beauville1990} represents multiplication by $\lambda$ as an
affine matrix function of $\nu$. Its leading matrix has eigenvalues
$1,-1,0$, with multiplicities $L,M,N$. The distinct branch intercepts
allow the three constant diagonal blocks to be diagonalized; the
remaining six blocks are precisely the unrestricted mode blocks
of \cref{eq:canonical-X}. \Cref{sec:geometry-converse}
gives the construction, proving~1.
\end{proof}

\subsection{Reconstructing a pencil from its spectral curve}
\label{sec:geometry-converse}

We prove the converse used in \cref{thm:geometry-classification}.
The construction follows the spectral line-bundle method
\cite[\S1]{Beauville1990}, with the prescribed infinity branches
determining the repeated eigenvalues of the leading matrix.

Let $C$ be the normalization of the plane model in statement~2 of
\cref{thm:geometry-classification}. Choose projective coordinates so
that the common line is $\tau=0$ and its three marked plane points
are $[1:1:0]$, $[-1:1:0]$, $[0:1:0]$.
Transversality makes $\nu$ have $d$ simple poles. Thus
\[
 \pi=\nu:C\longrightarrow\mathbb P^1,\qquad
 \pi^{-1}(\infty)=\mathcal D_\infty
 =\mathcal D_L+\mathcal D_M+\mathcal D_N
\]
is a finite degree-$d$ map, and the poles of $\lambda$ are bounded
by $\mathcal D_\infty$.

Choose a line bundle $\mathcal E$ of degree $G-1$ with
$H^0(C,\mathcal E)=0$, and put
$\mathcal L=\mathcal E(\mathcal D_\infty)$. Such $\mathcal E$
exist outside the theta divisor. The pushforward $\pi_*\mathcal L$
is a rank-$d$ bundle on $\mathbb P^1$ of degree zero, by
Riemann--Roch. Moreover, the projection formula gives
\[
 H^0\bigl(\mathbb P^1,(\pi_*\mathcal L)(-1)\bigr)
 =H^0(C,\mathcal E)=0.
\]
Every vector bundle on $\mathbb P^1$ splits into line bundles.
The vanishing makes all their degrees nonpositive, while their sum
is zero. Therefore $\pi_*\mathcal L\simeq\mathcal O^{d}$.

Multiplication by $\lambda$ now defines
\[
 \mathcal O^{d}\longrightarrow\mathcal O^{d}(1).
\]
In a trivialization it is an affine pencil $\mathsf X+\nu\mathsf D$.
At $\nu=\infty$ its leading operator is multiplication by
$\lambda/\nu$ on the $d$ separate fibers over infinity. Its values
are $1,-1,0$, with multiplicities $L,M,N$; hence a constant basis
change makes $\mathsf D=\diag(\Id_L,-\Id_M,0_N)$.
The characteristic polynomial is the plane equation: over a generic
$\nu$, the operator has the $d$ values of $\lambda$ as eigenvalues.

The eigenvalues of the three diagonal blocks of $\mathsf X$ are
the branch intercepts. Ordinary multiplicity gives distinct
intercepts within each block, so a block-diagonal basis change
diagonalizes all three blocks while preserving $\mathsf D$.
Write their entries as $-2b_\ell^L$, $-2b_m^M$, $b_n^N$,
with $1\leq\ell\leq L$, $1\leq m\leq M$, $1\leq n\leq N$.
The six remaining blocks are unrestricted. Assigning them the
factors and signs of \cref{eq:canonical-X} gives the three mode
matrices and their independent partners, and hence initial data
for \cref{eq:canonical-H}. This proves the required realization.


\subsection{Proof of the hyperelliptic criterion}

\begin{proof}[Proof of \cref{prop:geometry-hyperelliptic}]
Two singleton groups give a degree-two map in
\cref{eq:geometry-degrees}, proving sufficiency.
For necessity, sort the three sizes as $a\leq b\leq c$ and
suppose $b\geq2$. Two of the generating functions have degrees
$a+b$ and $a+c$, and
\[
 G-(a+c-1)=(a+b-2)c+(a-1)(b-2)>0.
\]
Both degrees are therefore at most $G$.
On a hyperelliptic curve of genus $G$, every meromorphic function
of degree at most $G$ factors through the degree-two hyperelliptic
map. Indeed, otherwise that map and the function would generate
the function field, and the Castelnuovo--Severi inequality would give
$G\leq\deg(f)-1$; see \cite[Theorem~14]{KhawajaSiksek2024}.
Both generating functions would thus belong to the rational
function field of the hyperelliptic map, a contradiction.
\end{proof}
